\documentclass{article}
\usepackage{graphicx}



\begin{document}
Estudo de Animais: Domésticos e Selvagens
Kawã Mendes
18 de Outubro de 2025

Estudo de Animais: Domésticos e Selvagens
Kawã Mendes
18 de Outubro de 2025

Seção: introdução

O objetivo deste estudo é apresentar informações básicas sobre animais domésticos e selvagens, classificando-os, como pode ser visto na Seção % ref classificação dos animais
, ou até expondo seus comportamentos e alimentação na Seção % ref Comportamentos e alimentação.
Incluímos figuras, tabelas e equações para ilustrar características e comportamentos de diferentes espécies.  
Referências são usadas para reforçar dados científicos % Citação lampert.

Animais domésticos, como cães e coelhos, convivem com seres humanos e possuem adaptações específicas para o ambiente doméstico.  
A Tabela % ref tabela animais
mostra exemplos de animais domésticos e suas características.

        Animal  Peso Médio (kg) Expectativa de Vida (anos)
        Cão  20  12 
        Gato  5  15 
        Coelho  2  8 

Subseção : Animais Selvagens

Animais selvagens vivem em seu habitat natural e não são domesticados.  
A Figura % ref figura elefante 
apresenta um exemplo de animal selvagem.

% < COLOQUE A FIGURA DO ELEFANTE AQUI >

Seção : Comportamento e alimentação

O comportamento dos animais pode ser descrito por fórmulas simples de energia.  
Por exemplo, a taxa metabólica E pode ser estimada por:

    E = k . m{0.75}

onde m é a massa do animal e k é uma constante específica da espécie.  
A Equação % ref equação
mostra como a energia diária depende do tamanho do animal.

Seção : Conclusão
Este documento apresentou um estudo introdutório sobre animais domésticos e selvagens, incluindo suas características, alimentação e comportamento.  
O uso de figuras, tabelas e equações facilita a compreensão dos conceitos.

\end{document}